% Modernised reading — fonds Alexandre Grothendieck, shelfmark 132
% (pages 1-31, the whole folder).
%
% This is an INTERPRETATION, not a transcription. It re-reads the pages in
% today's notation and vocabulary, states the hypotheses the manuscript leaves
% implicit, and drops the critical apparatus so that the mathematics reads
% straight through. Everything the brackets and underlines carried moves into a
% footnote. It is held to being mathematically correct as it stands; where the
% manuscript is loose or mistaken, the true statement is what appears here and
% the footnote says what the page has. In any dispute of substance the
% transcription governs: whoever cites Grothendieck must cite batch-01.fr.tex
% (pp. 1-20) or batch-02.fr.tex (pp. 21-31).
%
% Pass: Opus 5.5 (claude-opus-5-5), 2026-10-03 — first pass, unchecked against the pages by a human.
% The folder was read whole, its two batches together; this is one document
% for the shelfmark. It was made on Opus 5.5 and has not been measured against
% the reference reading of folder 115 (made on Opus 5). The literature named
% in the footnotes is cited from memory and was not checked against the
% sources.
%
% Scope. Pages 4, 9, 12 and 24 are blank. Pages 25-31 are two printed Notes
% aux Comptes rendus by Yves Ladegaillerie (1974); following RIGHTS.md and the
% transcription, they are identified and not restated: the reading says only
% what the transcription's notes say (title, volume, paragraphs).
%
% Spine. Three runs of manuscript, not one argument, plus the offprints.
%   pp. 1-3    fibre product of spaces vs 2-fibre product of fundamental
%              groupoids (complete as an argument)
%   pp. 5-11   groupoid models of surfaces with holes, the « Teichmüller à
%              trous » T and « spécial » TS, the sequence of page 8; the
%              combinatorial data of page 10 (ribbon graphs), page 11 breaks off
%   pp. 13-23  extensions of topological groups by a discrete N and base
%              change, his sections (A) p. 16, (B) p. 19, (B') pp. 20-21,
%              (C) pp. 22-23; page 18, interleaved, is a programme on
%              embeddings of 1-complexes in surfaces naming « Yves »
%   pp. 25-31  the offprints
%
% Notation fixed for the folder. Run 1 renames his G, H, f_1, g_1, phi to
% Gamma, Delta, f_*, g_*, theta (G, H and phi are used later for other
% things); his « Cat G » is written BG throughout. Run 2 keeps G, l_i, I, T,
% TS, phi : D -> C. Run 3 keeps N, G, H, pi, H~, H~_0, gothic G, z(N), phi;
% the boundary pi_1 H -> N of page 15 is the same phi as page 15's
% pi -> z(N) and is written phi from page 14 on. The symmetric group is
% \mathfrak{S} (page 18's letter, transcribed as gothic G, is read as S_n).
%
% Substantive departures, each footnoted where it occurs:
%   - p. 1: « connexes » read as path-connected (otherwise « f surjectif »
%     fails);
%   - p. 3: the pi_1 criterion is stated at the component of z_0; at another
%     component the condition is to be read there (images of pi_2 move by
%     the action of pi_1(S));
%   - pp. 5-8: « 2-spécial » is not defined; the reading takes (G, (l_i))
%     to be the fundamental group of a compact orientable surface with its
%     boundary loops, which is what « sphère à 1 trou » etc. presuppose;
%   - p. 7: last term « v_1(l_i) » read l_i;
%   - p. 13 onwards: standing hypotheses supplied (quotient maps are Serre
%     fibrations; groups locally path-connected and semi-locally simply
%     connected, e.g. Lie groups);
%   - p. 14: the square written with H/N^0, H'/N^0 is read with G/N^0,
%     G'/N^0, and the « isomorphism » is on pi_0, the « epi » on pi_1;
%   - p. 16: « trivialisation stricte » is given a definition (equivariant
%     splitting centralising N); the page's gloss is illegible;
%   - p. 17: Ext(gothic G, N) x Hom(pi, z(N)) is a product only once the
%     outer action of gothic G on N is fixed;
%   - p. 20: « extension centrale » (uncertain reading) of H by z(N) is
%     central only if gothic G acts trivially on z(N); « pi, N, H discrets »
%     read with gothic G for H;
%   - p. 21: the claim marked « faux » in the margin is replaced by what
%     holds: an exact sequence with kernel phi(pi)/phi(pi ∩ H~^0).
% Announced and not established: the classification of embeddings sketched
% on page 18; the category of page 11; the conclusion of (C) on page 23
% (struck); the boxed map of page 18.
\documentclass[11pt,a4paper]{article}
\input{../preamble/grothendieck}

\folder{132}
\pages{1}{31}
\dating{1974}
\watermark{Édition de démonstration\\interprétation personnelle de l'œuvre}
\foldertitle{Plongements dans surfaces (Ladegaillerie) : notes manuscrites (s.d.), tirés à part (1974) — lecture modernisée du dossier entier}

\begin{document}

\section*{Résumé}

\begin{resume}
Ces feuillets cherchent à décrire des surfaces — et les façons de dessiner un
réseau de fils sur une surface — par des objets purement algébriques, assez
fidèles pour qu'on puisse y lire leurs symétries.

Le point de départ est une idée qui traverse toute l'œuvre de Grothendieck :
un espace, vu « à homotopie près en dimension 1 », n'est rien d'autre qu'un
\emph{groupoïde}, c'est-à-dire l'ensemble de ses points et de ses chemins à
déformation près. Pour une surface à bord — une sphère percée de trous, un
tore troué —, ce groupoïde est très simple : un groupe libre, et dans ce
groupe, un lacet distingué pour chaque trou, qui fait le tour du bord. Le
dossier montre que cette donnée algébrique (un groupe et ses lacets de bord)
suffit à calculer un groupe qu'il appelle groupe de Teichmüller « à trous »,
et qui est, on le sait aujourd'hui, le groupe des symétries de la surface à
déformation près. Il trouve même deux
versions de ce groupe, selon qu'on autorise ou non les bords à tourner sur
eux-mêmes, et une suite exacte qui les relie : la différence, ce sont les
« tours de vis » qu'on peut donner à chaque bord, un entier par trou. Les
seules exceptions sont le disque et le cylindre, qui ont trop de symétries
continues pour se laisser décrire ainsi, et le dossier les calcule à part.

Pourquoi cela intéresse-t-il quelqu'un qui étudie des réseaux tracés sur une
surface ? Parce qu'un graphe dessiné sur une surface en a un voisinage — un
ruban épaissi, lui-même une surface à bord — et que la surface entière s'en
déduit en recollant, le long des bords, ce qui reste. Classer les dessins
revient alors à classer des recollements, et compter les symétries d'un
dessin revient à combiner les symétries du ruban et celles des morceaux
complémentaires. C'était le sujet des deux Notes aux Comptes rendus que Yves
Ladegaillerie, à Montpellier, publiait en 1974 et dont les tirés à part
ferment le dossier ; la page 18, la seule à le nommer, esquisse le programme
qui en découle et dit que le calcul du genre et du nombre de trous des
morceaux « est donné par Yves ».

Autour de ce noyau, deux études plus générales. La première (pages 1 à 3)
compare, pour un produit fibré d'espaces, le groupoïde du produit et le
produit des groupoïdes, et mesure exactement l'écart : il est nul dès que les
espaces n'ont pas de topologie au-delà de la dimension 1. La seconde (pages 13
à 23) est une petite théorie des extensions d'un groupe topologique par un
groupe discret : comment connaître les composantes connexes d'un groupe à
partir de son quotient, de son revêtement universel et d'une donnée finie, et
comment cette description se comporte quand on change de base. C'est la
question qu'on rencontre dès qu'on veut calculer le groupe des composantes
d'un groupe d'homéomorphismes, qui est précisément un groupe de Teichmüller ;
les pages ne font pas ce lien explicitement, mais elles posent les outils
qu'il demande.

Le dossier n'est pas un texte achevé : trois chantiers, ouverts côte à côte,
dont le premier seulement va au bout. On y voit une manière de travailler —
remplacer un objet géométrique par la plus petite structure algébrique qui le
retient, puis vérifier, cas par cas, ce que cette structure oublie ; ici, ce
qu'elle oublie se réduit à deux surfaces, le disque et le cylindre.

Les noms d'aujourd'hui : groupoïde fondamental, produit fibré homotopique,
groupe des classes d'applications (groupe modulaire) d'une surface à bord,
théorème de Dehn--Nielsen--Baer, twists de Dehn le long du bord, graphes
rubanés, revêtement universel d'un groupe topologique non connexe,
cohomologie non abélienne des extensions.

\keywords{fundamental groupoid, homotopy pullback, groupoid pullback,
Mayer-Vietoris sequence, surface with boundary, peripheral structure,
mapping class group, Dehn-Nielsen-Baer theorem, boundary Dehn twist,
ribbon graph, embeddings of graphs in surfaces,
extension of topological groups, universal covering group, Taylor
obstruction, non-abelian group extension, band (lien)}
\end{resume}

\pagerange{1}{31}
\section*{Le fil du dossier, et les conventions}

\subsection*{Les stations}

Le dossier compte trente et une pages : vingt de sa main, quatre blanches (4,
9, 12, 24), et sept pages imprimées (25 à 31), les tirés à part qui donnent au
dossier son titre. Les pages manuscrites ne sont pas datées ; la date de
l'inventaire, 1974, est celle des tirés à part.\footnote{Le dossier est rangé
dans le groupe « Descentes » de l'inventaire, daté « [avant 1970] » : la date
du groupe et celle du dossier ne s'accordent pas, et rien dans les pages ne
tranche.} Il n'y a pas un argument suivi, mais trois chantiers et un dépôt :

\begin{itemize}
\item \textbf{Produits fibrés et groupoïdes} — pages 1 à 3. Il porte sa
propre pagination, (1) à (3), et va au bout.
\item \textbf{Les surfaces à trous comme groupoïdes} — pages 5 à 11. Une
2-catégorie de « modèles » de surfaces à bord, le calcul de ses
automorphismes, deux groupes de Teichmüller et une suite exacte (pages 5 à
8) ; puis la donnée combinatoire d'un 1-complexe sur une surface et le
foncteur qui en fait un tel modèle (pages 10 et 11), qui s'interrompt.
\item \textbf{Extensions de groupes topologiques} — pages 13 à 23. Une
question (page 13), une réduction (pages 14 et 15), puis ses paragraphes
cerclés (A) page 16, (B) page 19, (B') pages 20 et 21, (C) pages 22 et
23.\footnote{Page 22, le « (C) » est écrit sur un « (D) » biffé.} La page
18, intercalée au milieu, appartient en réalité au deuxième chantier : c'est
un programme de classification des plongements de 1-complexes dans les
surfaces.
\item \textbf{Les tirés à part} — pages 25 à 31 : deux Notes d'Yves
Ladegaillerie, rangées dans l'ordre inverse de leur parution.
\end{itemize}

Ce que les trois chantiers ont en commun : chacun remplace un objet
topologique par sa donnée homotopique en basse dimension (un groupoïde, un
groupe et ses lacets, un groupe discret de composantes) et cherche ce que ce
remplacement perd.

\subsection*{Les conventions}

Les trois chantiers emploient les mêmes lettres pour des choses différentes ;
on les sépare ainsi.

\emph{Pages 1 à 3.} $f : X \to S$, $g : Y \to S$, $Z = X \times_S Y$ ;
$\Gamma = \pi_1(X, x_0)$, $\Delta = \pi_1(Y, y_0)$, $\Sigma = \pi_1(S, s_0)$,
et $f_* : \Gamma \to \Sigma$, $g_* : \Delta \to \Sigma$ les morphismes
induits.\footnote{Il écrit $G$, $H$, $\Sigma$ et $f_1$, $g_1$ ; les lettres
$G$ et $H$ sont réservées ici au troisième chantier.} $\Pi_1(X)$ est le
groupoïde fondamental, $\mathrm{B}\Gamma$ le groupoïde à un objet de groupe
$\Gamma$.\footnote{Il écrit $\mathrm{Cat}\,G$ ; on garde la notation
$\mathrm{B}G$ dans tout le document.}

\emph{Pages 5 à 11 et 18.} $G$ est le groupe fondamental d'une surface à
bord, $(\ell_i)_{i \in I}$ ses lacets de bord, $\varphi : \mathcal{D} \to
\mathcal{C}$ le foncteur « bord » ; $T$ et $TS$ ses deux groupes de
Teichmüller ; $\mathfrak{S}_I$ le groupe symétrique.

\emph{Pages 13 à 23.} $1 \to N \to G \to H \to 1$ est une extension de
groupes topologiques ; $\pi$, $\widetilde{H}$, $\widetilde{H}_0$,
$\mathfrak{G}$ sont les données du paragraphe (A) ; $\mathfrak{z}(N)$ est le
centre de $N$ ; $\varphi$ est l'homomorphisme $\pi \to \mathfrak{z}(N)$ qui
classe l'extension. Pour $K$ un groupe topologique, $K^0$ est la composante
neutre.

\emph{Hypothèses générales du troisième chantier.} Les pages ne les disent
pas ; elles sont nécessaires à presque chaque phrase. Les groupes
topologiques sont supposés localement connexes par arcs et semi-localement
simplement connexes — les groupes de Lie le sont —, et les projections
$K \to K/L$ sur un quotient par un sous-groupe fermé distingué sont des
fibrations de Serre ; $\pi_0$ désigne les composantes connexes par
arcs.\footnote{Hypothèses ajoutées par la lecture. Sans elles, ni la suite
exacte d'homotopie de la page 13, ni l'existence d'un revêtement universel à
la page 15, ni l'identification de $\pi_0$ à un quotient discret à la page 19
ne tiennent.}

\subsection*{Ce que le dossier annonce et n'établit pas}

\begin{itemize}
\item la catégorie des 1-complexes topologiques de la page 11 et le foncteur
de la page 10 sont posés, ni l'un ni l'autre n'est étudié ;
\item le programme de la page 18 — classer les plongements d'un 1-complexe
dans une surface et en calculer les groupes de symétries — s'arrête sur une
flèche encadrée « à déterminer » ;
\item l'affirmation de la page 21 sur $\pi_0(G)$ est marquée en marge d'un mot
lu « faux », et elle l'est en général ;
\item le paragraphe (C) des pages 22 et 23 établit le changement de base,
mais la phrase qui devait en tirer $\pi_0$ du groupe obtenu est barrée ; le
problème posé page 13 n'est donc pas résolu sur la page, bien que toutes les
pièces y soient (on dit lesquelles à la fin).
\end{itemize}

\pagerange{1}{2}
\section*{Produit fibré d'espaces et produit fibré de groupoïdes (pages 1 à 3)}

\subsection*{L'énoncé}

Soient $f : X \to S$ et $g : Y \to S$ deux applications continues, et
$Z = X \times_S Y$. Les deux projections et la commutativité du carré donnent
un foncteur
\[
  \theta : \Pi_1(Z) \longrightarrow
  \Pi_1(X) \overset{2}{\times}_{\Pi_1(S)} \Pi_1(Y)
\]
vers le \emph{produit fibré 2-catégorique} des groupoïdes : ses objets sont
les triplets $(x, y, \sigma)$ avec $\sigma$ une flèche de $f(x)$ vers $g(y)$
dans $\Pi_1(S)$, ses flèches les couples de flèches compatibles avec les
$\sigma$.\footnote{Il écrit le $2$ au-dessus du signe produit ; à sa seconde
occurrence, page 1, le signe se lit « ? » ou un 2 mal formé.}

\textbf{Proposition (page 1).} \emph{Si $f$ est une fibration de Serre,
$\theta$ est essentiellement surjectif et plein : il induit une bijection sur
les $\pi_0$ et, en chaque point, une surjection sur les groupes
fondamentaux.}

On peut supposer $X$, $Y$, $S$ connexes par arcs et non vides : chaque
composante connexe par arcs de $X$ se projette par une fibration de Serre
\emph{sur} une composante de $S$, puisqu'un chemin de $S$ se relève dès que
son origine est atteinte.\footnote{La page dit « connexes », et en tire « donc
$f$ surjectif ». C'est vrai pour la connexité par arcs, qui est d'ailleurs la
seule pertinente pour $\pi_0$ et $\pi_1$ ; pour la connexité au sens
topologique, une fibration de Serre au-dessus d'un espace connexe non connexe
par arcs n'est pas nécessairement surjective.} Choisissons $y_0 \in Y$,
$s_0 = g(y_0)$, $x_0 \in X$ au-dessus de $s_0$, et posons
$\Gamma$, $\Delta$, $\Sigma$ comme dans les conventions.

\subsection*{Le modèle algébrique}

À équivalence près, le produit fibré des groupoïdes est
$P = \mathrm{B}\Gamma \overset{2}{\times}_{\mathrm{B}\Sigma} \mathrm{B}\Delta$,
qui se décrit explicitement : ses objets sont les éléments $\sigma \in
\Sigma$, et une flèche $\sigma \to \sigma'$ est un couple $(a, b) \in \Gamma
\times \Delta$ tel que
\[
  \sigma' = g_*(b)\, \sigma\, f_*(a)^{-1} .
\]
Autrement dit, $\Gamma \times \Delta$ opère sur l'ensemble $\Sigma$ par
$(a, b) \cdot \sigma = g_*(b)\, \sigma\, f_*(a)^{-1}$, et $P$ est le
groupoïde de cette action. Ses classes d'isomorphie sont donc les doubles
classes $g_*(\Delta) \backslash \Sigma / f_*(\Gamma)$, et le groupe
d'automorphismes de $\sigma$ est son stabilisateur dans $\Gamma \times
\Delta$.\footnote{La page écrit les couples tantôt $(g, h)$, tantôt
$(h, g)$ ; on a uniformisé. Un petit carré en bas de page explicite la
condition de commutation.}

\subsection*{Les composantes}

La projection $Z \to Y$ est une fibration de Serre (image inverse de $f$), de
fibre $F = f^{-1}(s_0)$. Sa suite exacte d'homotopie donne
$\pi_0(Z) = \pi_0(F) / \pi_1(Y)$, l'ensemble des orbites de l'action de
$\pi_1(Y)$ sur $\pi_0(F)$ (puisque $Y$ est connexe par arcs).\footnote{À la
suite de la formule, une parenthèse se lit, avec doute, « coinvariants »,
suivie d'un mot illisible. C'est bien d'un ensemble d'orbites qu'il s'agit.}
De même, la fibration $F \subset X \to S$ donne $\pi_0(F) \simeq \Sigma /
f_*(\Gamma)$. D'où
\[
  \pi_0(Z) \simeq \Delta \backslash \Sigma / \Gamma ,
\]
qui est l'ensemble des classes de $P$ : $\theta$ est bijectif sur les
$\pi_0$.\footnote{La page commence ensuite (fin de la page 2) une preuve
directe de la surjectivité sur les $\pi_1$, par relèvement d'arcs, qu'elle
interrompt ; la page 3 la remplace par la comparaison des suites exactes
qu'on suit ici.}

\pagerange{3}{3}
\subsection*{Les groupes fondamentaux : l'écart exact}

Le groupoïde $\Phi$ des objets $\sigma \in \Sigma$, avec pour flèches
$\sigma \to \sigma'$ les $a \in \Gamma$ tels que $\sigma = \sigma' f_*(a)$, est
la fibre homotopique de $\mathrm{B}\Gamma \to \mathrm{B}\Sigma$ : son $\pi_0$
est $\Sigma / \Gamma$ et tous ses groupes d'automorphismes sont isomorphes à
$\operatorname{Ker} f_*$. Le foncteur $\Pi_1(F) \to \Phi$ est bijectif sur
les $\pi_0$ et surjectif sur les $\pi_1$, de noyau l'image de $\pi_2(S)$
dans $\pi_1(F)$, c'est-à-dire $\operatorname{Coker}(\pi_2(X) \to \pi_2(S))$ :
c'est la suite exacte d'homotopie de $F \subset X \to S$.

Au point $\sigma = 1$, le groupe d'automorphismes de $P$ est
\[
  \pi_1(P, 1) = \{ (a, b) \in \Gamma \times \Delta \mid f_*(a) = g_*(b) \} ,
\]
extension de $g_*^{-1}(f_*\Gamma)$ par $\operatorname{Ker} f_*$. Le groupe
$\pi_1(Z, z_0)$, $z_0 = (x_0, y_0)$, est extension du même groupe
$g_*^{-1}(f_*\Gamma)$ — l'image de $\pi_1(Z) \to \pi_1(Y)$ — par
$\operatorname{Coker}(\pi_2(Y) \to \pi_1(F))$. En comparant les deux :

\textbf{Proposition (page 3).} \emph{Le morphisme $\pi_1(Z, z_0) \to
\pi_1(P, 1)$ est surjectif, de noyau
\[
  \operatorname{Coker}\bigl(\pi_2(X, x_0) \times \pi_2(Y, y_0)
  \longrightarrow \pi_2(S, s_0)\bigr) ,
\]
le quotient de $\pi_2(S)$ par la somme des images de $\pi_2(X)$ et de
$\pi_2(Y)$. C'est un isomorphisme si et seulement si $\pi_2(X) \times
\pi_2(Y) \to \pi_2(S)$ est surjectif.}\footnote{La page fait le calcul au
point $\sigma = 1$, c'est-à-dire sur la composante de $z_0$, et énonce le
critère sans restriction. Sur une autre composante, celle de la double classe
de $\sigma$, il faut prendre les points base au-dessus d'un même point de
$S$, et les images de $\pi_2(X)$ et de $\pi_2(Y)$ s'y comparent à travers
l'action de $\sigma$ sur $\pi_2(S)$ : la condition est à vérifier composante
par composante.}

C'est, en termes d'aujourd'hui, la suite exacte de Mayer--Vietoris d'un
\emph{produit fibré homotopique} — $Z$ en est un, puisque $f$ est une
fibration — comparée à celle du produit fibré des groupoïdes, qui s'arrête
en degré 1.\footnote{Les fibrations de groupoïdes et le groupoïde fondamental
d'un produit fibré ont été étudiés par R. Brown (« Fibrations of groupoids »,
\emph{J. Algebra} 15, 1970, de mémoire) ; rien ne dit que ces pages le
connaissent, et rien ne permet de dater l'une par rapport à l'autre.}

Enfin, si $X$, $Y$, $S$ n'ont pas de $\pi_i$ pour $i \geqslant 2$, il en est
de même de $F$ (par $\pi_{i+1}(S) \to \pi_i(F) \to \pi_i(X)$), donc de $Z$ ;
la condition sur $\pi_2$ est alors vide, et $\theta$ est une équivalence :
\emph{pour des espaces de type $K(\pi, 1)$, le groupoïde du produit fibré est
le produit fibré des groupoïdes.}\footnote{La phrase finit sur un mot
illisible : « $X, Y, S$ sont des \dots ». « Asphériques », ou « $K(\pi,1)$ »,
conviendrait au sens ; la restitution est nôtre.}

\pagerange{5}{5}
\section*{Les surfaces à trous comme groupoïdes (pages 5 à 8)}

\subsection*{La 2-catégorie des modèles}

Un \emph{modèle} est un triplet $(\mathcal{C}, \mathcal{D}, \varphi)$ :
$\mathcal{D}$ un groupoïde dont chaque composante a un groupe fondamental
infini cyclique muni d'un générateur (une orientation),
$\mathcal{C}$ un groupoïde, $\varphi : \mathcal{D} \to \mathcal{C}$ un
foncteur. Le modèle à avoir en tête est
$\mathcal{C} = \Pi_1(\Sigma)$, $\mathcal{D} = \Pi_1(\partial\Sigma)$ pour une
surface compacte orientée $\Sigma$, les composantes du bord orientées comme
bord, et $\varphi$ induit par l'inclusion.\footnote{Il appelle $\mathcal{D}$
« groupoïde 1-spécial orienté », $\mathcal{C}$ « groupoïde 2-spécial »,
$\varphi$ « foncteur spécial », sans définir ces termes dans le dossier.
Tout ce qui suit — les exceptions « sphère à 1 trou », « sphère à 2 trous »,
le fait que $G$ soit libre — montre qu'il s'agit des groupoïdes fondamentaux
d'une surface compacte orientable à bord et de son bord ; c'est ainsi qu'on
lit la définition manquante.}

Une 1-flèche $(\mathcal{C}, \mathcal{D}, \varphi) \to (\mathcal{C}',
\mathcal{D}', \varphi')$ est un triplet $(u, v, \alpha)$ : une équivalence
$u : \mathcal{D} \to \mathcal{D}'$ compatible aux orientations, une
équivalence $v : \mathcal{C} \to \mathcal{C}'$, et un isomorphisme naturel
$\alpha : \varphi' u \Rightarrow v \varphi$. Une 2-flèche
$(u_1, v_1, \alpha_1) \Rightarrow (u_2, v_2, \alpha_2)$ est un couple
$(\lambda, \mu)$ d'isomorphismes naturels $\lambda : u_1 \Rightarrow u_2$,
$\mu : v_1 \Rightarrow v_2$, tel que le carré
\[
  \begin{array}{ccc}
    \varphi' u_1 & \xrightarrow{\ \varphi' \ast \lambda\ } & \varphi' u_2 \\
    {\scriptstyle \alpha_1} \big\downarrow & & \big\downarrow {\scriptstyle \alpha_2} \\
    v_1 \varphi & \xrightarrow{\ \mu \ast \varphi\ } & v_2 \varphi
  \end{array}
\]
commute.

\subsection*{La forme normale}

Si $\mathcal{C}$ est connexe et $\mathcal{D}$ non vide, le modèle est
équivalent à un modèle où $\mathcal{C} = \mathrm{B}G$ et
$\mathcal{D} = \coprod_{i \in I} \mathrm{B}\mathbb{Z}$ ; $\varphi$ est alors
donné par des éléments $\ell_i = \varphi_i(1) \in G$, les \emph{lacets de
bord}. Pour une surface de genre $g$ à $n = \operatorname{card} I \geqslant 1$
trous, $G$ est libre de rang $2g + n - 1$, présenté par
$\langle a_1, b_1, \dots, a_g, b_g, \ell_1, \dots, \ell_n \mid
[a_1, b_1] \cdots [a_g, b_g]\, \ell_1 \cdots \ell_n = 1 \rangle$.

Dans cette forme, et pour un second modèle $(G', (\ell'_{i'})_{i' \in I'})$,
une 1-flèche est la donnée
\begin{itemize}
\item d'une bijection $u : I \to I'$,
\item d'un isomorphisme $v : G \to G'$,\footnote{La page écrit
$v : G \xrightarrow{\sim} G$, sans accent au but.}
\item d'éléments $\alpha(i) \in G'$ tels que $v(\ell_i) = \alpha(i)\,
\ell'_{u(i)}\, \alpha(i)^{-1}$ ;
\end{itemize}
et une 2-flèche $(u, v_1, \alpha_1) \Rightarrow (u, v_2, \alpha_2)$ (il n'y
en a pas si $u_1 \neq u_2$) est un couple $(\lambda, \mu) \in \mathbb{Z}^I
\times G'$ tel que
\[
  v_2 = \operatorname{int}(\mu) \circ v_1 , \qquad
  \alpha_2(i)^{-1}\, \mu\, \alpha_1(i) = {\ell'_{u(i)}}^{\lambda_i}
  \quad (i \in I).
\]

\pagerange{6}{6}
\subsection*{Les composantes du groupe des automorphismes}

Prenons les deux modèles égaux. Les auto-équivalences forment un groupe
monoïdal (une \emph{Gr-catégorie}) $\underline{\mathrm{End}}$, dont on
calcule le $\pi_0$ et le $\pi_1$.

Si $(u_1, v_1, \alpha_1)$ et $(u_2, v_2, \alpha_2)$ sont isomorphes, alors
$u_1 = u_2 = u$ et $v_2 = \operatorname{int}(\mu) v_1$ pour un $\mu \in G$ ;
en écrivant $v_2(\ell_i)$ de deux façons, on trouve que
$\alpha_2(i)^{-1} \mu\, \alpha_1(i)$ centralise $\ell_{u(i)}$. Or, sauf pour
le disque, chaque lacet de bord engendre un sous-groupe infini cyclique qui
est son propre centralisateur.\footnote{Il dit « chaque $\ell_j$ est un
élément libre, et en tous cas son centralisateur est le groupe engendré par
$\ell_j$ ». Dans un groupe libre, le centralisateur d'un élément non trivial
est le sous-groupe cyclique maximal qui le contient ; il faut donc que
$\ell_j$ ne soit pas une puissance propre. C'est vrai : pour $n \geqslant 2$,
$\ell_j$ fait partie d'une base ; pour $n = 1$, $\ell_1$ est un produit de
commutateurs dont une puissance propre donnerait de la torsion dans le groupe
de la surface fermée. La justification est nôtre ; la page l'affirme.}
L'élément cherché est donc une puissance $\ell_{u(i)}^{\lambda_i}$, et
l'isomorphisme existe. D'où :

\textbf{Proposition (page 6).} \emph{$\pi_0(\underline{\mathrm{End}})$
s'identifie au sous-groupe de $\operatorname{Out}(G)$ formé des classes
$\overline{v}$ telles que, pour tout $i$, $\overline{v}(\ell_i)$ soit conjugué
à l'un des $\ell_j$.}

La permutation $u$ est en effet déterminée par $\overline{v}$, car deux lacets
de bord distincts ne sont jamais conjugués : leurs images dans l'abélianisé
$G_{\mathrm{ab}}$ sont déjà distinctes.\footnote{Il écrit
$\operatorname{Autext} G$ pour $\operatorname{Out}(G)$. Le $u$ du NB est
écrit comme un $\sigma$ sur la page.} Ce groupe est ce que la page appelle le
\emph{groupe de Teichmüller à trous} $T(\mathcal{C}, \mathcal{D}, \varphi)$.
On y reconnaît le groupe des classes d'isotopie d'homéomorphismes de la
surface préservant l'orientation, les composantes du bord pouvant être
permutées et l'isotopie n'étant pas astreinte à les fixer : c'est le
théorème de Dehn--Nielsen--Baer pour les surfaces à bord.\footnote{L'énoncé
topologique n'est pas sur la page, qui ne parle que de groupes. La condition
« conjugué à $\ell_j$ » et non à $\ell_j^{-1}$ traduit la préservation de
l'orientation. L'identification et le nom sont les nôtres (référence de
mémoire : Zieschang, Vogt, Coldewey, \emph{Surfaces and planar
discontinuous groups}, 1980) ; le nom « Teichmüller » est le sien.}

\pagerange{7}{7}
\subsection*{Le groupe fondamental}

$\pi_1(\underline{\mathrm{End}})$ est le groupe des automorphismes de
l'identité $(\mathrm{id}, \mathrm{id}, 1)$ : les $(\lambda, \mu)$ avec
$\mu \in Z(G)$, le centre, et $\mu = \ell_i^{\lambda_i}$ pour tout $i$.

\begin{itemize}
\item Hors du disque et du cylindre, $G$ est libre de rang au moins 2, de
centre trivial, et les $\ell_i$ sont d'ordre infini : $\mu = 1$,
$\lambda = 0$, et $\pi_1(\underline{\mathrm{End}}) = 0$.
\item Pour le cylindre (sphère à deux trous), $G \simeq \mathbb{Z}$,
$\ell_1 = 1$, $\ell_2 = -1$ en notation additive, $G$ est son propre centre,
et $\mu$ détermine $\lambda = (\mu, -\mu)$ : $\pi_1(\underline{\mathrm{End}})
= G \simeq \mathbb{Z}$, l'isomorphisme dépendant du choix d'un
générateur.\footnote{La page : « (pas canoniquement) ».}
\item Pour le disque (sphère à un trou), $G = 1$ et $\lambda \in \mathbb{Z}$
est libre : $\pi_1(\underline{\mathrm{End}}) \simeq \mathbb{Z}$,
canoniquement.
\end{itemize}

Les deux exceptions sont les deux surfaces dont le groupe d'homéomorphismes
a des composantes non simplement connexes : le disque et le cylindre tournent
sur eux-mêmes.

\subsection*{Le bord fixé : le groupe de Teichmüller spécial}

Fixons maintenant $\mathcal{D}$ : on ne garde que les 1-flèches avec
$u = \mathrm{id}$ et les 2-flèches avec $\lambda = 0$. On obtient une
Gr-catégorie $\underline{\mathrm{End}}_{\mathcal{D}}$ dont les objets sont
les couples $(v, \alpha)$, $v \in \operatorname{Aut} G$, $\alpha \in G^I$,
avec $v(\ell_i) = \alpha(i)\, \ell_i\, \alpha(i)^{-1}$, et dont les flèches
$(v_1, \alpha_1) \to (v_2, \alpha_2)$ sont les $\mu \in G$ tels que
$v_2 = \operatorname{int}(\mu) \circ v_1$ et $\alpha_2(i) = \mu\,
\alpha_1(i)$. Deux objets isomorphes ont même classe $\overline{v}$ dans
$\operatorname{Out}(G)$, et si $v_2 = \operatorname{int}(\mu_0) v_1$, alors
$\alpha_2(i)^{-1} \mu_0\, \alpha_1(i)$ centralise $\ell_i$.\footnote{Dans le
calcul de la page 7, le dernier terme est écrit
$\operatorname{int}(\mu_0)\, \operatorname{int} \alpha_1(i)\, v_1(\ell_i)$ ;
il faut $\ell_i$ au lieu de $v_1(\ell_i)$.} On appelle $TS$, groupe de
Teichmüller \emph{spécial}, le groupe $\pi_0(\underline{\mathrm{End}}_{
\mathcal{D}})$. Son $\pi_1$ est trivial dans tous les cas : un automorphisme
de $(\mathrm{id}, 1)$ est un $\mu$ avec $\mu \cdot 1 = 1$.

\pagerange{8}{8}
\subsection*{La suite exacte (page 8)}

\textbf{Proposition (page 8).} \emph{Il y a une suite
\[
  0 \longrightarrow \mathbb{Z}^I \longrightarrow TS
  \xrightarrow{\ \Psi\ } T \longrightarrow \mathfrak{S}_I \longrightarrow 1 ,
\]
où $\Psi(v, \alpha) = \overline{v}$, $T \to \mathfrak{S}_I$ envoie
$\overline{v}$ sur la permutation $u$ qu'il induit sur les lacets de bord,
et $\mathbb{Z}^I \to TS$ envoie $\lambda$ sur la classe de
$(\mathrm{id}, (\ell_i^{\lambda_i})_i)$. Elle est exacte, sauf en
$\mathbb{Z}^I$ pour le disque et le cylindre :
\begin{itemize}
\item pour le cylindre, $\operatorname{Ker} \Psi \simeq \mathbb{Z}$,
quotient de $\mathbb{Z}^2$ par l'antidiagonale, via la somme ;
\item pour le disque, $TS = 0$.
\end{itemize}}

L'exactitude en $T$ est immédiate : $\overline{v}$ fixe chaque classe de
conjugaison $[\ell_i]$ si et seulement si on peut choisir les $\alpha(i)$. La
surjectivité sur $\mathfrak{S}_I$ ne découle pas des calculs de la page ;
elle est vraie parce que toute permutation des composantes du bord d'une
surface connexe orientable se réalise par un homéomorphisme préservant
l'orientation.\footnote{La page écrit la suite sans justifier ce point ; la
justification topologique est nôtre.} Pour le noyau de $\Psi$ : un $(v,
\alpha)$ avec $\overline{v} = 1$ est isomorphe à un $(\mathrm{id}, \alpha)$,
et alors $\alpha(i)$ centralise $\ell_i$, donc $\alpha(i) =
\ell_i^{\lambda_i}$ avec $\lambda_i$ bien déterminé (hors du disque). Deux
tels objets $(\mathrm{id}, (\ell_i^{\lambda_i}))$ et $(\mathrm{id},
(\ell_i^{\lambda'_i}))$ sont isomorphes si et seulement s'il existe $\mu \in
Z(G)$ avec $\ell_i^{\lambda'_i} = \mu\, \ell_i^{\lambda_i}$ pour tout $i$.
Hors du cylindre, $Z(G) = 1$ et $\operatorname{Ker} \Psi \simeq
\mathbb{Z}^I$. Pour le cylindre, la condition s'écrit $\lambda'_1 =
\lambda_1 + \mu$, $\lambda'_2 = \lambda_2 - \mu$ : on identifie $\lambda$ et
$\lambda'$ quand leur différence est dans l'antidiagonale, et
$\operatorname{Ker} \Psi \simeq \mathbb{Z}$ par la somme.\footnote{En marge
de la suite : « exacte sauf pour la sphère à 1 ou 2 trous » (le premier mot
est de lecture incertaine). La page écrit, pour le cylindre,
$\operatorname{Ker}(\mathbb{Z}^I \to \operatorname{Ker} \Psi) \simeq
\operatorname{Ker}(\mathbb{Z}^I \to \mathbb{Z})$, « opération somme ».}

En termes d'aujourd'hui, $TS$ est le groupe des classes d'applications de la
surface relatif au bord (les homéomorphismes et les isotopies fixent le bord
point par point), $T$ celui où le bord est libre, et $\mathbb{Z}^I$ est
engendré par les twists de Dehn le long des composantes du bord ; la suite
est la suite bien connue qui relie le groupe relatif au bord à celui de la
surface où chaque bord est remplacé par une pointe, avec ses deux exceptions
classiques : sur le cylindre, les deux twists de bord coïncident ; sur le
disque, le twist de bord est isotope à l'identité relativement au bord
(lemme d'Alexander).\footnote{Ces identifications sont les nôtres ; la page
ne parle que des groupes $T$ et $TS$, qu'elle appelle « Teichmüller à trous »
et « Teichmüller spécial ». Que le $\pi_1$ de
$\underline{\mathrm{End}}_{\mathcal{D}}$ soit trivial dans tous les cas est
l'écho algébrique du fait que les composantes du groupe des homéomorphismes
fixant le bord sont contractiles.}

\pagerange{10}{10}
\section*{Le modèle d'un 1-complexe sur une surface (pages 10 et 11)}

La page 10 décrit les données combinatoires dont on fabrique un modèle. En
marge, de bas en haut le long de toute la moitié supérieure : « catégorie
\dots{} topologique des surfaces orientées ».\footnote{Un mot illisible après
« catégorie », et « orientées » de lecture incertaine.} Les données sont :
\begin{itemize}
\item[(a)] un 1-complexe combinatoire $(\Sigma, \omega)$ muni d'une
« circulation » ordinaire, c'est-à-dire d'un ordre cyclique des arêtes en
chaque sommet — ce qu'on appelle aujourd'hui un \emph{graphe rubané} —, sans
points isolés ni « bandes » isolées ;
\item[(b)] un groupoïde $\Pi_b$ à nombre fini de composantes, de groupes
fondamentaux infinis cycliques : les cercles isolés ;
\item[(c)] un tel groupoïde $\Pi_d$, orienté.
\end{itemize}
À ces données il associe :
\begin{itemize}
\item[(a$'$)] le 1-complexe orienté $\partial_t(\Sigma, \omega)$ des
« co-circuits » de $\Sigma$, et le foncteur $\varphi_0 :
\Pi_1(\partial_t\Sigma) \to \Pi_1(\Sigma)$ ;
\item[(b$'$)] pour chaque composante $\Pi_b(i)$, de groupe $\Lambda(i) \simeq
\mathbb{Z}$, le groupoïde $\widetilde{\Pi}_b(i) =
\coprod_{\mathrm{Isom}(\mathbb{Z}, \Lambda(i))} \Pi_b(i)$ à deux
composantes orientées, et $\varphi_b : \widetilde{\Pi}_b \to \Pi_b$ ;
\item[(c$'$)] $\varphi_d : \Pi_d \to \mathrm{B}\,\pi_0(\Pi_d)$, qui écrase
chaque composante sur un point ;
\end{itemize}
et le modèle
\[
  T\bigl((\Sigma, \omega), \Pi_d, \Pi_b\bigr) =
  \bigl(\Pi_1(\Sigma) \amalg \Pi_b \amalg \mathrm{B}\,\pi_0(\Pi_d),\
  \Pi_1(\partial_t\Sigma) \amalg \widetilde{\Pi}_b \amalg \Pi_d,\
  \varphi_0 \amalg \varphi_b \amalg \varphi_d\bigr) .
\]
Géométriquement, et c'est ainsi qu'on lit la page : le graphe rubané
s'épaissit en une surface à bord dont les composantes du bord sont ses
circuits de faces (les co-circuits) ; un cercle isolé s'épaissit en un
cylindre, dont les deux bords sont le cercle parcouru dans ses deux
orientations (d'où les deux composantes de $\widetilde{\Pi}_b(i)$, et le
cylindre de la page 7 avec $\ell_1 = 1$, $\ell_2 = -1$) ; un point isolé
s'épaissit en un disque, de groupoïde trivial, bordé par un cercle orienté.
$T$ associe ainsi à un 1-complexe le modèle de son \emph{voisinage
tubulaire}.\footnote{Cette lecture géométrique est la nôtre ; la page donne
les seules formules. Elle explique l'ordre des lettres (a), (c), (b) réunies
sur la page par une accolade, et le « NB $\varphi$ est foncteur bordant »
(dernier mot de lecture incertaine) qui suit la formule. Le mot qualifiant
$\widetilde{\Pi}_b(i)$ se lit, avec doute, « grossier ». Pour le genre et le
nombre de bords d'un graphe rubané connexe, le calcul est celui de la
caractéristique d'Euler : le nombre de bords est le nombre de circuits de
faces, et $2 - 2g - n = \#\text{sommets} - \#\text{arêtes}$.}

Sous un trait horizontal, la page commence la donnée d'un plongement : (d) un
groupoïde $\mathcal{C}'$, (e) un foncteur « bordant » $\varphi' : \mathcal{D}
\to \mathcal{C}'$, où $\mathcal{D}$ est celui de $T$. C'est, on le
comprendra page 18, le modèle du complémentaire, recollé au voisinage
tubulaire le long du bord.

\pagerange{11}{11}
La page 11 ouvre la catégorie, « non strictifiée », des 1-complexes
topologiques compacts « divisés » : des couples $(K_1, \Pi_b)$, $K_1$ un
1-complexe sans cercles isolés et $\Pi_b$ comme en (b), avec
\[
  \underline{\mathrm{Hom}}\bigl((K_1, \Pi_b), (K'_1, \Pi'_b)\bigr) \simeq
  \mathrm{Cat}\bigl(\mathrm{Isom}(K_1, K'_1)\bigr) \times
  \underline{\mathrm{Equ}}(\Pi_b, \Pi'_b) ,
\]
la catégorie discrète des isomorphismes combinatoires, fois le groupoïde des
équivalences. La formule court au bord de la feuille et le reste de la page
est blanc.\footnote{L'indice de $\underline{\mathrm{Equ}}$ est illisible ;
« combinatoires » est écrit sous la première ligne, et rattaché à la suivante
par une flèche.}

\pagerange{13}{13}
\section*{Extensions de groupes topologiques : la question (page 13)}

Soit une extension de groupes topologiques
\[
  (E) \qquad 1 \to N \to G \to H \to 1
\]
et un morphisme $u : H' \to H$ ; soit $G' = G \times_H H'$ l'extension de
$H'$ par $N$ obtenue par image inverse. La suite exacte d'homotopie de la
fibration $G \to H$ se termine par
\[
  (S) \qquad \pi_1 H \longrightarrow \pi_0 N \longrightarrow \pi_0 G
  \longrightarrow \pi_0 H \longrightarrow 1 ,
\]
de même $(S')$ pour $G'$, et $u$ induit un morphisme $(S') \to (S)$. Le carré
de droite donne
\[
  c : \pi_0 G' \longrightarrow \pi_0 G \times_{\pi_0 H} \pi_0 H' ,
\]
toujours surjectif (on relève dans $G$ un chemin de $H$ joignant deux points
de même composante), mais pas toujours injectif : la seule connaissance de
$(S)$ et de $u$ en degrés $0$ et $1$ ne détermine pas $\pi_0 G'$.

\emph{Le problème} : trouver des données algébriques simples attachées à
$(E)$ qui, jointes à $\pi_0 u$ et $\pi_1 u$, déterminent $\pi_0 G'$ à
isomorphisme près.\footnote{La page écrit « connaissant $\pi_0 H \to \pi_0 H'$
et $\pi_1 H \to \pi_1 H'$ » : les flèches sont dans le sens inverse de $u$.
Elle ajoute au-dessus de la ligne « (discrets) », de lecture et de point
d'attache incertains. « Algébriques » est de lecture incertaine.}

\pagerange{14}{14}
\section*{Réduction, et le cas d'une base connexe (pages 14 et 15)}

\subsection*{$\alpha$) Réduction au cas $N$ discret}

La composante neutre $N^0$ est caractéristique dans $N$, donc distinguée dans
$G$ et dans $G'$. Remplaçons $G$ par $G/N^0$, extension de $H$ par $N/N^0 =
\pi_0 N$ ; l'image inverse de $G/N^0$ par $u$ est $G'/N^0$, et la base $H$,
$H'$, $u$ ne change pas. Comme $N^0$ est connexe, $\pi_0 G \to \pi_0(G/N^0)$
et $\pi_0 G' \to \pi_0(G'/N^0)$ sont bijectifs, et $\pi_1 G \to
\pi_1(G/N^0)$ est surjectif.\footnote{La page dessine ce carré avec
$H/N^0$ et $H'/N^0$, où l'on attend $G/N^0$ et $G'/N^0$ ; elle dit les
flèches horizontales « isom. » sur les $\pi_i$ puis « épi. » sur les $\pi_1$.
Ce qui est vrai, et suffit à la réduction, est l'énoncé donné ici : bijection
sur $\pi_0$, surjection sur $\pi_1$. Elle écrit aussi « $\pi_0 G / N^0$ »
pour $\pi_0(G/N^0)$, et « extension de $H$ » pour $G'/N^0$, qui est une
extension de $H'$.} Le problème ne change pas : \emph{on peut supposer $N$
discret.}

Alors $G \to H$ est un revêtement, $(S)$ se lit
\[
  \pi_i G \simeq \pi_i H \ (i \geqslant 2), \qquad
  \pi_1 G = \operatorname{Ker}(\varphi), \qquad
  1 \to \operatorname{Coker}(\varphi) \to \pi_0 G \to \pi_0 H \to 1 ,
\]
où $\varphi : \pi_1 H \to N$ est le bord de la suite d'homotopie. Son image
est contenue dans le centre $\mathfrak{z}(N)$ : si $(g_t)$ est un chemin de
$G$ issu de $1$, la conjugaison par $g_t$ sur le groupe discret $N$ est
constante, donc l'extrémité d'un relèvement de lacet commute à
$N$.\footnote{Remarque nôtre ; elle explique pourquoi, page 15 et au-delà,
$\varphi$ est à valeurs dans le centre, et pourquoi $\operatorname{Coker}
\varphi$ est un groupe.} Pour $G'$, le bord est $\varphi' = \varphi \circ
\pi_1 u$, connu dès que $\varphi$ et $\pi_1 u$ le sont. Il reste à
déterminer l'extension $\pi_0 G'$ de $\pi_0 H'$ par $R' =
\operatorname{Coker}(\varphi')$.

\pagerange{15}{15}
\subsection*{Le noyau de $c$}

L'image inverse de l'extension $\pi_0 G$ par $\pi_0 H' \to \pi_0 H$ est
$\pi_0 G \times_{\pi_0 H} \pi_0 H'$, extension de $\pi_0 H'$ par $R =
\operatorname{Coker}(\varphi)$ — et non par $R'$. Or $R$ est le quotient de
$R'$ par $\operatorname{Im} \varphi / \operatorname{Im} \varphi'$, d'où la
suite exacte
\[
  1 \longrightarrow \frac{\pi_1 H}{\operatorname{Ker} \varphi +
  \operatorname{Im} \pi_1(u)} \longrightarrow \pi_0 G' \xrightarrow{\ c\ }
  \pi_0 G \times_{\pi_0 H} \pi_0 H' \longrightarrow 1 .
\]
Elle dit exactement comment le défaut d'injectivité de $c$ dépend de
$\pi_1 H$ et de $\pi_1 u$ : c'est un quotient de $\operatorname{Coker}
\pi_1(u)$.\footnote{La première phrase de la page 15 est de lecture
incertaine (« Montrons que », « non pas $R'$ », « $R$ est un quotient de
$R'$ ») ; le sens est fixé par la formule, qui est nette.}

\subsection*{$\beta$) Cas $H$ connexe}

Soit $\pi = \pi_1(H)$ et $\widetilde{H}$ le revêtement universel de $H$,
extension de $H$ par $\pi$ ; $\pi$ est discret et distingué dans le groupe
connexe $\widetilde{H}$, donc central. Alors :

\textbf{Proposition (page 15).} \emph{Si $H$ est connexe et $N$ discret, les
extensions de $H$ par $N$ sont classées par les morphismes $\varphi : \pi \to
\mathfrak{z}(N)$, et chacune n'a que l'identité pour automorphisme
d'extension. L'extension associée à $\varphi$ est l'image directe
$\widetilde{H} \times^{\pi} N = (\widetilde{H} \times N) / \{ (\gamma,
\varphi(\gamma)^{-1}) \}$.}

En effet, l'application $\widetilde{H} \to H$ se relève en un morphisme
$\widetilde{H} \to G$, dont l'image, connexe, centralise $N$ ; d'où
$\widetilde{H} \times N \to G$, surjectif, de noyau le graphe de
$\varphi^{-1}$ sur $\pi$. Un automorphisme d'extension est donné par une
application continue de $H$ connexe dans $N$ discret, donc
triviale.\footnote{La page dit l'extension « canonique » (lecture
incertaine), « à isom. unique près » ; la démonstration esquissée ici est
nôtre.} Pour $u : H' \to H$ quelconque — et c'est la remarque utile, car
$H'$ n'a pas à être connexe —, $G'$ s'obtient en tirant d'abord
$\widetilde{H}$ en $H' \times_H \widetilde{H}$, extension de $H'$ par $\pi$,
puis en poussant par $\pi \to \mathfrak{z}(N) \subset N$.\footnote{La fin de
la page est très incertaine (« Il n'y a pas de problème si $H'$ n'est pas
connexe », « décrivons les ext. de $\pi_0 H'$ \dots ») et se termine par deux
lignes biffées.}

\pagerange{16}{16}
\section*{(A) Le cadre général : base non connexe (pages 16 et 17)}

Quand $H$ n'est pas connexe, il n'a pas de revêtement universel au sens
habituel. On se donne donc à la place :
\begin{itemize}
\item[(a)] l'extension $1 \to N \to G \to H \to 1$, $N$ discret ;
\item[(b)] une extension de groupes topologiques $1 \to \pi \to
\widetilde{H} \to H \to 1$, $\pi$ discret ;
\item[(c)] un sous-groupe distingué $\widetilde{H}_0$ de $\widetilde{H}$ de
quotient discret $\mathfrak{G}$ : $1 \to \widetilde{H}_0 \to \widetilde{H}
\to \mathfrak{G} \to 1$.
\end{itemize}
On étudie, (b) et (c) fixés, les extensions $G$ de $H$ par $N$ \emph{munies}
d'une trivialisation stricte de leur image inverse sur $\widetilde{H}_0$,
c'est-à-dire d'un morphisme $\widetilde{H}_0 \to G$ qui relève
$\widetilde{H}_0 \to H$ et centralise $N$.\footnote{On appelle
\emph{trivialisation stricte} d'une extension $E$ de $K$ par $N$ sur un
sous-groupe distingué $L \subset K$ un morphisme $L \to E$ relevant
l'inclusion, centralisant $N$, et dont l'image est distinguée dans $E$
(équivariance pour la conjugaison). La page glose « stricte » par une
parenthèse ajoutée en interligne, lue avec un fort doute « comm. avec prod.
produits » ; on donne la définition dont la suite a besoin.} Si
$\widetilde{H}_0$ est connexe et simplement connexe, une telle
trivialisation existe et est unique : c'est la structure de revêtement.

Soit $\widetilde{G} = G \times_H \widetilde{H}$, extension de $\widetilde{H}$
par $N$ ; elle contient $\pi$ comme sous-groupe distingué centralisant $N$, et
$G = \widetilde{G}/\pi$. Se donner $G$ revient donc à se donner
$\widetilde{G}$ muni d'une trivialisation stricte sur $\pi$ ; la structure
supplémentaire de $G$ revient à une trivialisation stricte de $\widetilde{G}$
sur $\widetilde{H}_0$. Or une extension de $\widetilde{H}$ par $N$
strictement trivialisée sur $\widetilde{H}_0$ est la même chose qu'une
extension
\[
  1 \longrightarrow N \longrightarrow \widetilde{G}_0 \longrightarrow
  \mathfrak{G} \longrightarrow 1
\]
(on retrouve $\widetilde{G} = \widetilde{H} \times_{\mathfrak{G}}
\widetilde{G}_0$, et $\widetilde{G}_0 = \widetilde{G} / \widetilde{H}_0$).

\pagerange{17}{17}
Il reste la trivialisation sur $\pi$. Dans ces termes, c'est un morphisme
$\varphi : \pi \to \widetilde{G}_0$ qui relève le composé $\pi \subset
\widetilde{H} \to \mathfrak{G}$, centralise $N$, et est équivariant pour les
actions de conjugaison.\footnote{Le début de la page 17 (« i.e. un
relèvement, et compatible \dots ») est surchargé et se lit mal ; la flèche
$\varphi$ y est pointillée.} \emph{Cas particulier} : si $\pi \subset
\widetilde{H}_0$, le composé $\pi \to \mathfrak{G}$ est trivial, et
$\varphi$ est un morphisme
\[
  \varphi : \pi \longrightarrow \mathfrak{z}(N) ,
\]
équivariant pour l'action de $\mathfrak{G}$ — sur $\pi$ par conjugaison dans
$\widetilde{H}$, sur $\mathfrak{z}(N)$ par conjugaison dans
$\widetilde{G}_0$. D'où, dans ce cas, la classification :
\[
  \bigl\{ G \text{ muni de sa trivialisation sur } \widetilde{H}_0 \bigr\}
  \ \simeq\ \bigl\{ (\widetilde{G}_0, \varphi) \bigr\} ,
\]
les couples formés d'une extension $\widetilde{G}_0$ de $\mathfrak{G}$ par
$N$ et d'un $\mathfrak{G}$-morphisme $\varphi : \pi \to
\mathfrak{z}(N)$.\footnote{La page écrit
$\operatorname{Ext}(\mathfrak{G}, N) \times
\operatorname{Hom}_{\mathfrak{G}}(\pi, \mathfrak{z}(N)) \xrightarrow{\sim}
\operatorname{Ext}(H \bmod \widetilde{H}_0, N) \to \operatorname{Ext}(H, N)$.
Ce n'est un produit qu'une fois fixée l'action extérieure
$\mathfrak{G} \to \operatorname{Out}(N)$, car c'est elle qui définit l'action
de $\mathfrak{G}$ sur $\mathfrak{z}(N)$ ; sinon, c'est un ensemble de
couples fibré au-dessus de $\operatorname{Ext}(\mathfrak{G}, N)$.}

Suivent, sous un trait, des brouillons barrés de longues obliques (les suites
de $\widetilde{H}$ et d'un $\widetilde{H}'$ au-dessus de $H'$, une grille de
sous-groupes), puis trois lignes que la même oblique traverse sans qu'il soit
sûr qu'elle les annule : l'ensemble des extensions est un « espace homogène
sous $H^2(H, \mathfrak{z}(N))$ », avec le morphisme $H^2(\mathfrak{G},
\mathfrak{z}N) \to H^2(H, \mathfrak{z}N)$ induit par $H \to \mathfrak{G}$.
C'est la théorie des extensions non abéliennes : à action extérieure fixée,
les extensions d'un groupe par $N$, s'il en existe, forment un espace
principal homogène sous le $H^2$ à coefficients dans le centre, et
l'obstruction à leur existence est dans le $H^3$.\footnote{La page écrit
aussi « $1 \to N \to E \to \widetilde{H} \to 1$ équivaut à
$H^2(\mathfrak{G}, N)$ », qui ne se comprend qu'à ce sens près : un
torseur sous $H^2(\mathfrak{G}, \mathfrak{z}(N))$, et non $H^2(\mathfrak{G},
N)$. Pour $H$ topologique, le $H^2$ s'entend en un sens convenable
(cohomologie continue ou de Segal--Moore) ; la page ne précise pas.
Références de mémoire : Eilenberg--MacLane (1947) pour les groupes
discrets ; Giraud, \emph{Cohomologie non abélienne} (1971), pour la notion de
lien.}

\pagerange{18}{18}
\section*{Un programme : plonger un 1-complexe dans une surface (page 18)}

Intercalée entre (A) et (B), la page 18 revient aux surfaces. Elle est
organisée autour d'un schéma de sigles :

\begin{tikzcd}[column sep=small, row sep=small]
& \mathrm{SO} & & & \\
& & \mathrm{PLO} \arrow[ul] \arrow[dl] \arrow[drr] & & \\
& \mathrm{TUBO} \arrow[dl] \arrow[r] & \mathrm{SOB} \arrow[dr, "\partial"] & & \mathrm{SOB} \arrow[dl, "-\partial"'] \\
\mathrm{Kc} & & & \mathrm{CO} &
\end{tikzcd}

Les sigles ne sont pas développés.\footnote{Lecture nôtre, guidée par les
tirés à part du dossier, dont la seconde Note porte sur les « catégories TUB
et TUBO » de voisinages tubulaires et sur les « 1-complexes circularisés ».
Sous PLO : « (avec $\pi_0(K) \to \pi_0(S)$ surjectif) » ; sous CO, un
« $\simeq \mathrm{Equiv}(\mathbb{N}, \{\mathfrak{G}_n\})$ » dont le premier
mot est incertain, et que l'on comprend comme le groupoïde des ensembles finis,
somme des $\mathrm{B}\mathfrak{S}_n$ (la lettre transcrite $\mathfrak{G}_n$
est lue ici comme le groupe symétrique, le $\mathfrak{S}_I$ de la page 8).
La flèche de droite vers CO porte « $-\partial$ ($\simeq \partial$) ».} On
les lit ainsi : PLO, les plongements d'un 1-complexe $K$ dans une surface
orientée ; SO, les surfaces orientées ; TUBO, les voisinages tubulaires
orientés ; Kc, les 1-complexes circularisés (graphes rubanés) ; SOB, les
surfaces orientées à bord ; CO, les collections finies de cercles orientés.
Le schéma dit alors qu'un plongement, c'est la surface ambiante, le voisinage
tubulaire du complexe (qui détermine le graphe rubané, et réciproquement à
équivalence près), et le complémentaire du voisinage, recollés le long du bord
où les deux orientations s'opposent : le $\partial$ et le $-\partial$ du
schéma. C'est le complément, en termes de modèles, de la page 10.

Autour du schéma, les annotations décomposent le problème des symétries :
\begin{itemize}
\item les \emph{complexes hypercirculaires} : un complexe circularisé $K$ avec
une partition de l'ensemble $\mathrm{Circ}(K)$ de ses circuits et un genre
$g \in \mathbb{N}$ par bloc — chaque bloc de circuits borde une composante
du complémentaire, de genre $g_i$ à $n_i$ trous ;
\item le groupe des symétries d'un plongement est une extension du groupe des
symétries du complexe hypercirculaire par les groupes de Teichmüller des
composantes du complémentaire, et cette extension est connue dès qu'on
connaît ces groupes séparément, les morphismes canoniques $\mathbb{Z}^{n_i}
\to$ (centre), et les extensions $1 \to T_{g_i, n_i} \to L_{g_i, n_i} \to
\mathfrak{S}_{n_i} \to 1$ ;\footnote{Les lettres de cette ligne ne
s'accordent pas exactement avec celles de la page 8 : $L_{g,n}$, qui se
projette sur le groupe symétrique et reçoit $\mathbb{Z}^{n}$ dans son centre,
tient à la fois du $T$ et du $TS$ de la page 8, et un passage encadré,
entièrement hachuré, où l'on devine « $L_{g_i, n_i} \to T_{g_i, n_i}$ de
noyau $\operatorname{Im}(\mathbb{Z}^{n_i} \to L_{g_i, n_i})$ », a été
abandonné. On ne tranche pas. Le premier groupe est noté $\Pi_{g_i, n_i}$,
« spéciaux » étant de lecture incertaine.}
\item tout se ramène aux groupes de Teichmüller « à trous » ordinaires
$T_{g,n}$ et aux groupes de symétries des 1-complexes circularisés connexes.
\end{itemize}
« Le calcul des $g_i n_i$ est donné par Yves » : le genre et le nombre de
trous des morceaux.\footnote{Le prénom se lit nettement ; c'est le seul lien
explicite des pages manuscrites avec Ladegaillerie.} La page se termine sur
ce qu'il reste à déterminer, encadré :
\[
  \operatorname{Aut}_{\mathrm{Circ}}(K) \longrightarrow
  \operatorname{Aut}\bigl(T_{?}(K)\bigr) \simeq T_{g,n} ,
\]
l'action des symétries d'un complexe circularisé sur le groupe de Teichmüller
de son épaississement.\footnote{L'indice de $T$ est illisible.} Le dossier
ne le détermine pas.

\pagerange{19}{19}
\section*{(B) Le cas du revêtement universel (pages 19 et 20)}

Supposons à la fois
\begin{itemize}
\item[$\alpha$)] $\widetilde{H}_0$ connexe et simplement connexe ;
\item[$\beta$)] $\pi \subset \widetilde{H}_0$ ;
\item[$\gamma$)] $\pi$, $N$, $\mathfrak{G}$ discrets.
\end{itemize}
Comme $\mathfrak{G}$ est discret, $\widetilde{H}_0$ est ouvert ; connexe, il
est la composante neutre $\widetilde{H}^0$, et $\mathfrak{G} = \pi_0
\widetilde{H}$. Comme $\pi \subset \widetilde{H}_0$, $H = \widetilde{H}/\pi$
est extension de $\mathfrak{G} = \pi_0(H)$ par $H^0 = \widetilde{H}_0/\pi$ ;
et comme $\pi$ est discret et $\widetilde{H}_0$ simplement connexe,
$\widetilde{H}_0$ est le revêtement universel de $H^0$, avec $\pi =
\pi_1(H^0) = \pi_1(H)$. La situation équivaut donc à la donnée de :

\begin{itemize}
\item[1°)] un groupe topologique $H$, extension du groupe discret
$\mathfrak{G} = \pi_0(H)$ par le groupe connexe $H^0$, de revêtement
universel $\widetilde{H^0}$, extension de $H^0$ par $\pi = \pi_1(H^0)$ ;
\item[2°)] une extension $\widetilde{H}$ de $H$ par $\pi$ qui prolonge
$\widetilde{H^0}$ — un « revêtement universel » du groupe non connexe $H$ ;
\item[3°)] le groupe discret $N$.
\end{itemize}

Une telle extension $\widetilde{H}$ n'existe pas toujours et n'est pas unique.
Si elle existe, les autres s'en déduisent par l'action de
$H^2(\mathfrak{G}, \pi)$, et l'obstruction à son existence est dans
$H^3(\mathfrak{G}, \pi)$ — « sauf erreur », dit la page. Elle existe si $H$
est un produit semi-direct $\mathfrak{G} \ltimes H^0$ : on prend
$\widetilde{H} = \mathfrak{G} \ltimes \widetilde{H^0}$, l'action de
$\mathfrak{G}$ se relevant au revêtement universel.\footnote{Le résultat est
connu : R. L. Taylor, « Covering groups of nonconnected topological groups »,
\emph{Proc. AMS} 5 (1954), et R. Brown, O. Mucuk, \emph{Math. Proc.
Cambridge Philos. Soc.} 115 (1994), qui situent l'obstruction dans
$H^3(\pi_0 H, \pi_1 H^0)$ ; références de mémoire, et rien ne dit que la
page connaît la première. La page dit l'extension « unique (i.e. mod. action
de $H^2(\mathfrak{G}, \pi)$) comme extension de $\mathfrak{G}$ par $\pi$ » :
on écrit ce qui tient, un espace homogène sous $H^2(\mathfrak{G}, \pi)$.}

\textbf{Proposition (page 19).} \emph{Sous $\alpha$), $\beta$), $\gamma$), la
catégorie des extensions de $H$ par $N$ est équivalente à celle des couples
$(\widetilde{G}_0, \varphi)$, où $\widetilde{G}_0$ est une extension de
$\mathfrak{G}$ par $N$ et $\varphi : \pi \to \mathfrak{z}(N)$ un
$\mathfrak{G}$-morphisme.}

C'est le cas particulier de (A), puisque sous $\alpha$) la trivialisation
sur $\widetilde{H}_0$ existe et est unique. On a $\widetilde{G}_0 = \pi_0
\widetilde{G}$, la composante neutre de $\widetilde{G}$ s'envoyant
isomorphiquement sur $\widetilde{H}^0$.\footnote{NB écrit en biais dans
l'angle inférieur gauche, de lecture partielle.} Pour $H$ connexe, on
retrouve la proposition de la page 15.

\pagerange{20}{20}
\subsection*{La construction explicite, et $\pi_0(G)$}

Réciproquement, $G$ se construit ainsi à partir de $(\widetilde{G}_0,
\varphi)$. Soit $G_!$ l'image inverse de $\widetilde{G}_0$ par $H \to
\mathfrak{G}$, extension de $H$ par $N$. Le $\mathfrak{G}$-morphisme
$\varphi$ pousse l'extension $\widetilde{H}$ de $H$ par $\pi$ en une
extension $\widetilde{H}^{\varphi}$ de $H$ par $\mathfrak{z}(N)$, où $H$ opère
sur $\mathfrak{z}(N)$ à travers $\mathfrak{G}$ comme il le fait par
l'intermédiaire de $G_!$.\footnote{La page dit l'extension « centrale »
(lecture incertaine) ; elle ne l'est que si $\mathfrak{G}$ opère
trivialement sur $\mathfrak{z}(N)$, et la page précise elle-même l'action.
L'indice de $G_!$ est un point d'exclamation, surmonté peut-être d'un
$\varphi$.} Le produit contracté de $G_!$ et de $\widetilde{H}^{\varphi}$
— le quotient de $G_! \times_H \widetilde{H}^{\varphi}$ par l'antidiagonale de
$\mathfrak{z}(N)$ — est le $G$ cherché : c'est l'action de
$H^2(H, \mathfrak{z}(N))$ sur les extensions de lien fixé de la page 17.

Comme $G$ est un revêtement de $H$, $\pi_i(G) \simeq \pi_i(H)$ pour $i
\geqslant 2$ et $\pi_1(G) = \operatorname{Ker}(\pi_1 H \to N)$, où le bord
$\pi_1(H) = \pi \to N$ est $\varphi$ suivi de $\mathfrak{z}(N) \subset N$.
Donc $\pi_0 G$ est extension de $\mathfrak{G}$ par
$\operatorname{Coker}(\varphi)$, et plus précisément
\[
  \pi_0(G) \simeq \operatorname{Coker}(\pi \to \widetilde{G}_0) =
  \widetilde{G}_0 / \varphi(\pi) ,
\]
$\varphi(\pi)$ étant distingué dans $\widetilde{G}_0$ par
équivariance.\footnote{La page écrit $G_0$ pour $\widetilde{G}_0$ dans cette
formule et dans la définition de $G_!$.} C'est la réponse à la question de
la page 13 pour un tel $G$.

\pagerange{20}{21}
\section*{(B$'$) Sans l'hypothèse $\alpha$) (pages 20 et 21)}

Supposons seulement $\beta$) et $\gamma$) : $\pi$, $N$, $\mathfrak{G}$
discrets, $\pi \subset \widetilde{H}_0$, mais $\widetilde{H}_0$ pas
nécessairement connexe ni simplement connexe.\footnote{La page écrit
« $\pi, N, H$ discrets », le $H$ appuyé ; la condition $\gamma$) de la page
19 porte sur $\mathfrak{G}$, et c'est ce qu'on lit.} Ce cas compte, parce que
c'est celui qu'on obtient par changement de base, au paragraphe (C). $N$ étant
discret, on a encore $\pi_i(G) \simeq \pi_i(H)$ pour $i \geqslant 2$, et
\[
  \pi_1(G) = \operatorname{Ker}\bigl(\pi_1(H) \to N\bigr) ,
\]
où le bord se factorise par $\pi_1(H) \xrightarrow{\partial} \pi_0(\pi) = \pi
\xrightarrow{\varphi} \mathfrak{z}(N) \subset N$, $\partial$ étant le bord de
la suite d'homotopie de $1 \to \pi \to \widetilde{H} \to H \to 1$. La suite
$\pi_1 H \to N \to \pi_0 G \to \pi_0 H \to 1$ fait de $\pi_0 G$ une extension
de $\pi_0 H$ par $\operatorname{Coker}(\varphi \partial)$.

L'image de $\partial$ est $\pi \cap \widetilde{H}^0$ (suite exacte
$\pi_1(H) \to \pi \to \pi_0(\widetilde{H})$). L'extension
$\widetilde{G}_0$ de $\mathfrak{G}$ par $N$ est encore définie, puisque
$\pi \subset \widetilde{H}_0$ ; par $\pi_0 H \to \mathfrak{G}$ elle donne
une extension $G_1 = \pi_0 H \times_{\mathfrak{G}} \widetilde{G}_0$ de
$\pi_0 H$ par $N$. La page affirme ensuite $\pi_0 G \simeq
\operatorname{Coker}(\pi_1 H \to G_1)$, et marque la phrase en marge d'un mot
lu « faux », suivi d'une croix. Ce qui tient est ceci :

\textbf{Proposition (lecture, page 21).} \emph{Le morphisme $G \to
\widetilde{G}_0/\varphi(\pi)$ induit une suite exacte
\[
  1 \longrightarrow \varphi(\pi) / \varphi(\pi \cap \widetilde{H}^0)
  \longrightarrow \pi_0 G \longrightarrow G_1 / \varphi(\pi)
  \longrightarrow 1 .
\]
En particulier $\pi_0 G \simeq G_1 / \varphi(\pi)$ si et seulement si
$\varphi(\pi) = \varphi(\pi \cap \widetilde{H}^0)$, ce qui a lieu sous
$\alpha$).}

En effet, $G = \widetilde{G}/\pi$ et $\widetilde{G} \to \widetilde{G}_0$
envoie $\pi$ sur $\varphi(\pi)$ ; d'où $G \to \widetilde{G}_0/\varphi(\pi)$,
compatible à $G \to H \to \mathfrak{G}$, qui se factorise par $\pi_0 G$ et
donne $\pi_0 G \to \pi_0 H \times_{\mathfrak{G}}
(\widetilde{G}_0/\varphi(\pi)) = G_1/\varphi(\pi)$, surjectif ; et les
noyaux sur $\pi_0 H$ sont $N/\varphi(\pi \cap \widetilde{H}^0)$ et
$N/\varphi(\pi)$.\footnote{Cette proposition est nôtre : elle remplace
l'affirmation barrée en marge, qui n'est exacte que dans le cas où le noyau
est trivial. Elle a la même forme que la suite exacte de la page 15, et pour
la même raison : le quotient par l'image du $\pi_1$ de la base ne voit pas
tout $\pi$.}

\pagerange{22}{23}
\section*{(C) Le changement de base (pages 22 et 23)}

Plaçons-nous sous les données (A), et soit $u : H' \to H$ un morphisme de
groupes topologiques. Par image inverse on obtient $\widetilde{H}' =
\widetilde{H} \times_H H'$, extension $(b')$ de $H'$ par le même $\pi$ ; on
note $\widetilde{H}'_0$ l'image inverse de $\widetilde{H}_0$ et
$\mathfrak{G}' = \widetilde{H}'/\widetilde{H}'_0$, qui s'injecte dans
$\mathfrak{G}$ ; c'est la suite $(a')$.\footnote{La page étiquette « (c) » la
ligne du bas, là où l'appariement avec (a$'$) ferait attendre « (a) ».}
L'extension $G' = G \times_H H'$ hérite d'une trivialisation stricte sur
$\widetilde{H}'_0$, image inverse de celle de $G$ sur $\widetilde{H}_0$. Par
(A), $G'$ est donc déterminée par
\begin{itemize}
\item[1°)] une extension $\widetilde{G}'_0$ de $\mathfrak{G}'$ par $N$ ;
\item[2°)] un morphisme $\varphi' : \pi \to \widetilde{G}'_0$ centralisant
$N$ et relevant $\pi \to \mathfrak{G}'$ — donc, si $\pi \subset
\widetilde{H}'_0$, un morphisme $\pi \to \mathfrak{z}(N)$.\footnote{Dans la
suite exacte de 1°), l'objet du milieu porte deux accents, là où le texte
n'en porte qu'un.}
\end{itemize}

\textbf{Proposition (page 23).} \emph{$\widetilde{G}'_0$ est l'image inverse
de $\widetilde{G}_0$ par $\mathfrak{G}' \hookrightarrow \mathfrak{G}$ (c'est
donc un sous-groupe de $\widetilde{G}_0$), et $\varphi'$ est induit par
$\varphi$ ; si $\pi \subset \widetilde{H}_0$, $\varphi' = \varphi$ comme
morphisme $\pi \to \mathfrak{z}(N)$.}

En effet $\widetilde{G}' = \widetilde{H}' \times_{\mathfrak{G}}
\widetilde{G}_0$, et son quotient par $\widetilde{H}'_0$ est
$\mathfrak{G}' \times_{\mathfrak{G}} \widetilde{G}_0$. Les conditions
$\beta$) et $\gamma$) passent aux données $(a')$, $(b')$, $(c')$ — $\pi
\subset \widetilde{H}'_0$ équivaut à $\pi \subset \widetilde{H}_0$ —, mais
pas $\alpha$) : $\widetilde{H}'_0$ n'a aucune raison d'être connexe. D'où le
paragraphe (B$'$). La phrase suivante, qui commençait à en tirer $\pi_0$, est
barrée de quatre longs traits.\footnote{Elle se lit en partie : « Mais comme
$\mathfrak{G}'$ est discret, il en est de même de $\mathfrak{G}'_0$, donc
$G'_0 \dots$ comme un quotient de $\pi_0 \widetilde{G}'$ ».}

\subsection*{Ce que les pièces donnent}

Le dossier ne l'écrit pas, mais les pièces suffisent à répondre à la question
de la page 13, sous $\alpha$), $\beta$), $\gamma$) pour $H$. Par (C),
$G'$ correspond au couple $(\widetilde{G}_0 \times_{\mathfrak{G}}
\mathfrak{G}', \varphi)$ ; par la proposition de la page 21 appliquée à
$G'$, avec $G'_1 = \pi_0 H' \times_{\mathfrak{G}} \widetilde{G}_0$,
\[
  1 \longrightarrow \varphi(\pi) / \varphi(\pi \cap \widetilde{H}'^{\,0})
  \longrightarrow \pi_0 G' \longrightarrow G'_1 / \varphi(\pi)
  \longrightarrow 1 ,
\]
et $\pi \cap \widetilde{H}'^{\,0}$ est l'image de $\pi_1(H') \to \pi_1(H) =
\pi$. Les données sont donc : l'extension $\widetilde{G}_0$ de $\mathfrak{G}$
par $N$, le $\mathfrak{G}$-morphisme $\varphi$, et $u$ en degrés $0$ et $1$ —
ce que la page 13 demandait. Le noyau de gauche est un quotient de
$\operatorname{Coker} \pi_1(u)$, comme à la page 15.\footnote{Ce paragraphe
est entièrement nôtre ; il assemble les propositions des pages 19, 21 et 23,
et n'a pas d'équivalent sur les pages. Sur le sens qu'aurait cette théorie
pour les groupes d'homéomorphismes des surfaces, dont le $\pi_0$ est un
groupe de Teichmüller, les pages ne disent rien.}

\pagerange{25}{31}
\section*{Les tirés à part (pages 25 à 31)}

Le dossier se termine par deux Notes imprimées d'Yves Ladegaillerie,
« Classification topologique des plongements des 1-complexes compacts dans
les surfaces », présentées par Henri Cartan aux \emph{Comptes rendus de
l'Académie des sciences de Paris}, série A. Les pages 25 à 28 sont la seconde
Note (t. 279, 22 juillet 1974, p. 129-132, §§ 7 à 13 : voisinages tubulaires
et catégories TUB et TUBO, 1-complexes circularisés, cocircularisations,
classification des plongements, plongements minimaux et genre de $K$) ; les
pages 29 à 31 la première (t. 278, 27 mai 1974, p. 1401-1403, §§ 1 à 5 :
catégorie des 1-complexes à isotopie près, préliminaires, voisinage tubulaire
d'un 1-complexe circularisé et d'un plongement, théorème de classification).
Elles sont rangées dans l'ordre inverse de leur parution. On n'y trouve rien
de sa main.\footnote{Ce sont des textes d'un autre auteur ; on se borne à
les identifier, comme la transcription. Un croquis au crayon au bas de la
page 28, sans rapport avec les mathématiques, n'est pas attribuable.} Le
vocabulaire des pages 10 et 18 — complexes circularisés, voisinages
tubulaires orientés, circuits — est celui de ces Notes ; les pages
manuscrites ne disent pas si elles les précèdent ou les suivent.

\end{document}
